
All the results of the \emph{Bible} (\emph{loc.\ cit.}) concerning semisimple
groups hold more generally for reductive groups (defined below), and their
proofs are identical after the following harmless modifications:
\begin{itemize}
\item Exposé 9, §4, Definition 3: see 1.6.1 below;
\item Exposé 12, §4, Theorem 2(e): omit ``finite'';
\item Exposé 13, §3, Theorem 2: replace ``rank'' by ``semisimple rank'';
\item Exposé 13, §4, Corollary 2 to Theorem 3: replace ``rank'' by
  ``reductive rank''.
\end{itemize}

\noindent\SGAInlineRefTarget{sga3:xix:paragraph:1-2}\textbf{1.2.}\quad
Let \(G\) be a smooth affine connected \(k\)-group.  The \emph{radical} of
\(G\) (\emph{Bible}, §9.4, Proposition 2)%
is the reduced subgroup associated with the identity component of the
intersection of the Borel subgroups of \(G\); it is also the largest smooth
connected solvable normal subgroup of \(G\).  We denote it by
\(\operatorname{rad}(G)\).

The \emph{unipotent radical} of \(G\) is the unipotent part of the radical of
\(G\); it is also the largest smooth connected unipotent normal subgroup of
\(G\).  We denote it by \(\operatorname{rad}^{u}(G)\).

\noindent\SGAInlineRefTarget{sga3:xix:paragraph:1-3}\textbf{1.3.}\quad
Let \(G\) be a smooth affine connected \(k\)-group, and let \(Q\) be a torus
in \(G\).  Then the centralizer
\(\underline{\operatorname{Centr}}_{G}(Q)\) of \(Q\) in \(G\) is a closed
subgroup of \(G\) (\SGARefLink{sga3:viii:corollary:6:7}{Exposé VIII, 6.7}), smooth (\SGARefLink{sga3:xi:corollary:2:4}{Exposé XI, 2.4}), and connected
(cf.\ \emph{Bible}, §6.6, Theorem 6(a), or [Ch05], §6.7, Theorem 7(a)).

We have the fundamental relation
\begin{equation}
  \operatorname{rad}^{u}\!\bigl(\underline{\operatorname{Centr}}_{G}(Q)\bigr)
  =
  \operatorname{rad}^{u}(G)\cap
  \underline{\operatorname{Centr}}_{G}(Q).
  \tag{1.3.1}\label{sga3:xix:eq:1:3:1}
\end{equation}

First,%
\footnote{Editor's note: the argument that follows has been expanded.}
\[
  U=\operatorname{rad}^{u}(G)\cap
  \underline{\operatorname{Centr}}_{G}(Q)
\]
is a normal unipotent subgroup of
\(\underline{\operatorname{Centr}}_{G}(Q)\); let us show that it is smooth
and connected.  Let \(Q\) act on \(\operatorname{rad}^{u}(G)\) by inner
automorphisms.  Then \(U\) is precisely the scheme of fixed points of this
action.  This scheme is smooth and connected by \SGARefLink{sga3:xix:lemma:1:4}{Lemma 1.4} below.

Consequently, \(U\) is a closed subgroup of
\(\operatorname{rad}^{u}(\underline{\operatorname{Centr}}_{G}(Q))\).
On the other hand, by the \emph{Bible}, Exposé 12, §3, Corollary to
Theorem 1,
\[
  \operatorname{rad}^{u}\!\bigl(
    \underline{\operatorname{Centr}}_{G}(Q)\bigr)(k)
  \subset
  \operatorname{rad}^{u}(G)(k).
\]
Equality \eqref{sga3:xix:eq:1:3:1} follows.

Let us record one special case of the preceding result.  If \(G\) is a smooth
affine connected \(k\)-group and \(T\) is a maximal torus of \(G\), then
\begin{equation}
  \underline{\operatorname{Centr}}_{G}(T)
  =
  T\cdot\bigl(
    \underline{\operatorname{Centr}}_{G}(T)\cap
    \operatorname{rad}^{u}(G)\bigr).
  \tag{1.3.2}\label{sga3:xix:eq:1:3:2}
\end{equation}
Indeed (cf.\ \emph{Bible}, §6.6, Theorem 6(c), or [Ch05], §6.7,
Theorem 7(c)), \(\underline{\operatorname{Centr}}_{G}(T)\) is solvable and
hence is the semidirect product of its maximal torus \(T\) and its unipotent
radical.

By the density theorem (cf.\ \emph{Bible}, §6.5, Theorem 5(a), or [Ch05],
§6.6, Theorem 6(a)), the union of
\(\underline{\operatorname{Centr}}_{G}(T)\), as \(T\) ranges over the
maximal tori of \(G\), contains a dense open subset of \(G\).  It therefore
follows from \eqref{sga3:xix:eq:1:3:2} that:
